\section{Introduction}\label{sec:introduction}

For a polynomial $h\in\Z[x]$, consider how large a set of integers can be
if no difference between two of its elements is a nonzero value of $h$.
Write $\N_+=\{1,2,\ldots\}$, $[N]=\{1,\ldots,N\}$, and
$A-A=\{a-b:a,b\in A\}$. We call $h$ \emph{intersective} if, for every
integer $q\ge1$, some integer $t$ satisfies $h(t)\equiv0\pmod q$.
This is the necessary local condition for polynomial differences in every
set of positive density: if $h$ has no root modulo $q$, the multiples of
$q$ contain no difference in $h(\N_+)$. The condition allows the roots at
different moduli to arise from different factors of $h$; a rational root is
not required.

\begin{theorem}\label{thm:main}
For every integer $k\ge2$, there is a constant $c_k\in(0,1)$ with the
following property. For every intersective polynomial $h\in\Z[x]$ of degree
$k$ with positive leading coefficient, there are a constant $C_h\ge1$ and an
integer $N_h\ge1$ such that, whenever $N\ge N_h$ and $A\subseteq[N]$ satisfy
\[
 (A-A)\cap h(\N_+)\subseteq\{0\},
\]
one has
\[
 |A|\le C_h N^{1-c_k}.
\]
\end{theorem}

The exponent is common to all polynomials of a fixed degree. The constants
$C_h,N_h$ may depend on every coefficient of $h$, and $h$ remains fixed as
$N$ grows. The theorem concerns ordinary integer differences; zero values of
$h$ impose no restriction on $A$.

\subsection{History and relation to earlier work}

Furstenberg~\cite{Furstenberg1977} and S\'ark\"ozy~\cite{Sarkozy1978}
proved that every set of positive upper density contains two elements
whose difference is a nonzero square. Their proofs introduced distinct
ergodic and Fourier-analytic approaches to this recurrence problem.
Kamae and Mend\`es France~\cite{KamaeMendesFrance1978} extended the
positive-density conclusion to nonconstant intersective polynomials. Together with
the elementary congruence obstruction above, this identifies the exact
class of nonconstant polynomials for which the density of an avoiding set must tend
to zero.

The quantitative problem has driven successive refinements of Fourier
density-increment methods. For square differences, Pintz, Steiger, and
Szemer\'edi~\cite{PintzSteigerSzemeredi1988} introduced a substantially
stronger iteration. Bloom and Maynard~\cite{BloomMaynard2022} later
obtained the bound
\[
 |A|\ll N(\log N)^{-c\log\log\log N}
\]
using estimates for the additive energy of rational numbers.
For general intersective polynomials, Lucier~\cite{Lucier2006} established
quantitative bounds using auxiliary polynomials adapted to progressions,
and Rice~\cite{Rice2019} developed stronger estimates by combining
iteration with sieved polynomial sums.
Arala~\cite{Arala2024} extended the Bloom--Maynard bound to every
intersective polynomial of degree at least two, with a constant depending
only on the degree in the exponent of $\log N$.

Green and Sawhney~\cite[Theorem~1.1]{GreenSawhney2025} obtained the
square-difference bound
\[
 |A|\ll N\exp(-c\sqrt{\log N}).
\]
Their arithmetic
level-$d$ inequality and random sparsification also underpin the recent
general-polynomial result of Adajar et al.~\cite[Theorem~1.1]{AdajarEtAl2026}:
for every fixed intersective $h$ of degree at least two and every fixed
$0<\mu<1/2$,
\[
 |A|\le N\exp\bigl(-c_{h,\mu}(\log N)^\mu\bigr)
\]
for sufficiently large $N$. Theorem~\ref{thm:main} gives a fixed power
of $N$ as the saving, with an exponent common to all polynomials of each
fixed degree.

The immediate predecessor for our method is the power-saving theorem
for square differences in~\cite[Theorem~1.1]{squarepower}. Its proof
uses finite reflection-positive tuple laws and a signed multilinear
functional. Hatami~\cite{Hatami2010} gave a multilinear H\"older criterion
for graph norms. Reflection and Cauchy--Schwarz arguments are central to
the reflection-group constructions of Conlon and Lee~\cite{ConlonLee2017}.
Here we adapt the specific laws and functional of~\cite{squarepower},
with all required proofs included, and supply the polynomial and
change-of-scale arguments described next.

\subsection{The argument and its principal constructions}

The square-difference method constructs a multilinear form
which dominates a fixed power of density, and proves that this form contracts
on passing to shorter intervals. We retain this organization, but polynomial
values require two additional compatibilities: the polynomial changes when
one passes to a progression, and the local distribution supplied by the
multilinear construction need not be the natural distribution of polynomial
values.

\paragraph{Normalizing the local coefficients.}
We first pass from the polynomial in Theorem~\ref{thm:main} to a polynomial
$g(x)=h(r+D_0x)/M$ for which the first nonzero Taylor coefficient at one
chosen root in each $\Z_p$ is a unit. Only finitely many primes require the
extra division by $M$. The multiplicity of each chosen root is retained, and
the representative $r$ is chosen so that positive inputs for $g$ give positive
inputs for $h$. The purpose is to make the local numerical choices that
determine the final exponent depend only on the degree.

\paragraph{A family closed under passage to progressions.}
Section~\ref{sec:local} retains the auxiliary construction for an arbitrary
fixed intersective polynomial $h$. It gives integer polynomials $h_l$, indexed by
$l\in\N_+$, with $h_1=h$, and a fixed completely multiplicative function
$\lambda:\N_+\to\N_+$ such that, for all positive integers $l,D$,
\[
 h_l(Dx-i)=\lambda(D)h_{lD}(x)
 \quad\text{for some }0\le i<D,
 \qquad D\mid\lambda(D),\quad \lambda(D)\le D^{\deg h}.
\]
Consequently, a subset of an arithmetic progression of step $\lambda(D)$
which avoids nonzero $h_l$-differences becomes, in progression coordinates,
a set avoiding nonzero $h_{lD}$-differences. Compatible roots in the rings
$\Z_p$ produce the family. Its coefficient bounds, complete exponential-sum
estimates, and local lifting properties are uniform in $l$.

Applied to $g$, the family retains the zero valuation of the first nonzero
coefficient at each chosen root: Lucier's auxiliary formula preserves that
valuation \cite[Section~8]{Lucier2006}. Thus every auxiliary polynomial for
$g$ is nonconstant as an element of $\F_p[x]$ for every prime $p$. For
$p>k$, its reduced derivative is nonzero, and the complete-sum estimates hold
beyond a threshold depending only on $k$. At the finitely many smaller primes below that
threshold, interpolation bounds the derivative valuation at some argument in
terms of $k$ and $p$. These bounds make the lifting moduli and cycle length,
and hence the number of tuple slots below, depend only on $k$.

\paragraph{Weighted forms that follow the family.}
For each sufficiently large prime $p$, Section~\ref{sec:tuples} constructs a
probability distribution on tuples $(z_v)_{v\in V}\in\F_p^V$, where $V$ is
one fixed finite set of even cardinality $d$. A distinguished cycle of slots
has increments among the values of $h_l$ at arguments where its derivative is
nonzero modulo $p$, apart from a small exceptional part of the distribution.
The constraint on an edge is weighted by the number of such arguments giving
its increment. This weight has a Fourier bound uniform in the auxiliary
polynomial and transforms exactly under the progression identity above.

The laws are reflection positive: the quadratic form obtained by splitting
the slots across any of a specified family of reflections is positive
semidefinite. They also have uniform individual marginals and a conditional
mixing estimate. Section~\ref{sec:lifts} derives the resulting signed
H\"older inequality and estimates for Fourier lifts. At scale $N$, the lift
retains the normalized Fourier coefficients of $1_A$ at rational frequencies
whose denominators are at most a small fixed power of $N$. Integrating its
product over the $d$ slots gives a nonnegative quantity $Y_l(N,A)$ satisfying
\[
 Y_l(N,A)\ge (|A|/N)^d.
\]
The associated seminorm permits comparison along progressions with a factor
tending to one, rather than a fixed loss at every scale.

\paragraph{A kernel realizing the resulting increment distribution.}
The increment distribution of two slots is determined by the whole tuple
construction. It is supported on regular polynomial values, but there is no
reason for it to equal the distribution obtained by evaluating $h_l$ at a
uniform argument. Section~\ref{sec:kernel} constructs a signed kernel on
actual positive values of $h_l$ whose Fourier transform approximates the
desired local multipliers. At each sufficiently large prime, a density on regular arguments
realizes the prescribed increment distribution. A truncated divisor
expansion combines these densities, while uniform lifting handles
characters of higher prime-power conductor. Minor arcs are controlled by
elementary Weyl differencing.

Only support is needed to use this kernel: its bilinear form vanishes on two
ordered subsets of an avoiding set, regardless of the signs of the kernel
weights. Its Fourier approximation therefore gives cancellation for the
chosen pair of slots. This separates the construction of a useful tuple law
from the problem of realizing its pair distribution on integer polynomial
values.

\paragraph{Contraction while the auxiliary parameter grows.}
Section~\ref{sec:transfer} combines the kernel estimate with the seminorm
comparison. A fixed positive proportion of assignments of residue classes
at the finitely many small primes gives cancellation. The other assignments,
together with the summable exceptional parts of the tuple laws, are bounded
using shorter intervals and polynomials farther along the family.
Section~\ref{sec:induction} closes the resulting recurrence with a bound
that holds for some fixed $a,b>0$, simultaneously for every $l$ and every set $A$
avoiding nonzero $h_l$-differences:
\[
 Y_l(N,A)\le C l^bN^{-a}.
\]
The factor $l^b$ is necessary to accommodate the
growth of the auxiliary polynomial: when $l$ is large relative to $N$, a
uniform subpower bound suffices, while the kernel supplies contraction in
the remaining range. For the normalized polynomial $g$, the parameters can
be chosen so that $a/d$ depends only on $k$. Taking $l=1$ and using density
domination bounds sets avoiding nonzero $g$-differences. Partitioning $A$
into residue classes modulo $M$, and using $Mg(n)=h(r+D_0n)$ with
$r+D_0n\ge1$, then transfers the bound to $h$ and proves
Theorem~\ref{thm:main}.

The polynomial argument thus adds exact covariance of the weighted tuple laws
across the auxiliary family and a signed polynomial kernel realizing their
pair multipliers. The initial normalization controls the structural parameters
that determine the exponent. All the estimates needed from the
reflection-positive method are proved below.

\subsection{Notation and uniformity}

We write $e(x)=\exp(2\pi i x)$ and use the negative sign in Fourier
coefficients. For $\xi\in\Q/\Z$, let $q(\xi)$ be its reduced positive
denominator, with $q(0)=1$. Averages, measures, and norms on finite sets use
uniform probability measure unless a different measure is specified.
For a positive integer $n$, write $\rad(n)=\prod_{p\mid n}p$, with
$\rad(1)=1$.
Empty products equal $1$, and a product space with no coordinates is a
singleton. Conditional expectations are with respect to the law specified
in the surrounding statement.

The auxiliary-family development applies to any fixed intersective polynomial
$h$ with chosen roots, and we write $k=\deg h$. Constants in $O(\cdot)$ and
$\ll$ may depend on this fixed datum and on parameters already fixed from it;
they are independent of the varying integers, sets, and auxiliary index $l$,
unless a restriction is explicitly stated. Likewise, sufficiently-large
thresholds are uniform in those varying data. We explicitly identify the
structural choices that depend only on $k$ when the first nonzero coefficient
at every chosen root is a unit. Bounds denoted $O_k(\cdot)$ are uniform over
all fixed data of degree $k$ satisfying those choices. A bound $N^{o(1)}$
means $C_\epsilon N^\epsilon$ for every $\epsilon>0$ under the same
uniformity convention. All structural parameters are fixed before the
asymptotic scale $N$ varies.
